% Paper 1 (two-way nondeterministic complementation): abstract, introduction lines 1-52, automata.tex lines 171-180
\begin{abstract}
We prove that two-way nondeterministic finite automata cannot be
complemented with a polynomial number of states independent of the
alphabet. For each $n\ge4$ we construct an $n$-state automaton over a finite
alphabet whose complement requires at least
$\tfrac12 2^{\lfloor(n-4)/127\rfloor}-1$ states.
\section{Introduction}
\label{sec:introduction}

A two-way nondeterministic finite automaton (2NFA) has a finite set of
states and a single read-only head, which may move in either direction
along an input bounded by endmarkers. It accepts when some finite
computation reaches an accepting state. The complementation problem asks
whether every $n$-state 2NFA over a finite alphabet $\Sigma$ has a 2NFA for
the complementary language with at most $p(n)$ states, for one polynomial
$p$ independent of $\Sigma$. We answer this question negatively.

For exact state counts, there is one initial state and a set of accepting
states. Each transition depends only on the current state and scanned
symbol, updates the state, and moves the head left, right, or not at all.
We start the head on the left endmarker and forbid moves beyond either
endmarker. Missing transitions
reject that branch; infinite nonaccepting computations do not accept.
An accepting initial configuration counts as acceptance. All states,
including initial and accepting states, are counted. A deterministic
two-way automaton (2DFA) has at most one successor for each state and
scanned symbol.

For a finite set $H$, let $\mathcal R_H$ denote all binary relations on
$H$. Products are composed in path order:
$(p,r)\in AB$ if and only if $(p,q)\in A$ and $(q,r)\in B$ for some $q$.
Use $\Sigma_H=\mathcal R_H$ as an input alphabet. The product $r(w)$ of a
word $w=R_1\cdots R_k$ is $R_1\cdots R_k$, and the empty product is
$r(\varepsilon)=I_H=\{(p,p):p\in H\}$. Define the relation-product
liveness language
\[
 L_H=\{w\in\Sigma_H^*:r(w)\ne\varnothing\}.
\]
In particular, $L_H$ contains the empty word when $H$ is nonempty.
For nonempty words, this is Sakoda and Sipser's family $B_h$, where
$h=|H|$ \cite[Section~2.1 and Theorem~2.3]{SakodaSipser1978}; the name
\emph{one-way liveness} is used in \cite[Section~2.2]{Kapoutsis2006}.

\begin{theorem}\label{main:theorem}
For every integer $n\ge4$, set $H=\{1,\ldots,n-2\}$ and
$\Sigma_n=\mathcal R_H$. There is an $n$-state 2NFA $A_n$ recognizing
$L_H$ such that every 2NFA recognizing
$\Sigma_n^*\setminus L(A_n)$ has at least
\[
 \frac12\,2^{\lfloor(n-4)/127\rfloor}-1
\]
states. In particular, there is no polynomial state bound for
complementation that is independent of the finite alphabet.
\end{theorem}

The alphabet $\Sigma_n$ consists of all binary relations on a set of
$n-2$ points, so $|\Sigma_n|=2^{(n-2)^2}$. The parameter is the
number of states, not the size of the transition table. Theorem~\ref{main:theorem}
\begin{corollary}\label{auto:determinization}
There are absolute constants $c>0$ and $n_0$ such that, for every
$n\ge n_0$, every 2DFA recognizing the language of the $n$-state 2NFA
$A_n$ in Theorem~\ref{main:theorem} has at least
\[
 c\,2^{\lfloor(n-4)/127\rfloor}
\]
states. Hence no polynomial state bound independent of the finite
alphabet can determinize all 2NFAs.
\end{corollary}

% Paper 2 (one-way liveness): abstract, introduction lines 64-130
\begin{abstract}
One-way liveness on $h$ points accepts a word of binary relations when
their ordered product is nonempty. For every $h\ge2$, it has a
nondeterministic automaton with $h+3$ states and no left moves, whereas
every equivalent $s$-state two-way deterministic automaton satisfies
$4(s+2)^2\ge2^{\lfloor(h-2)/31\rfloor}$. Partial transition rules, stay
moves, and nonaccepting infinite computations are allowed. The alphabets are
finite and grow with $h$, so the result rules out an alphabet-independent
polynomial state bound for deterministic two-way simulation, already for
one-way nondeterministic sources.
\end{abstract}
\subsection{Model and statement}

An automaton has a finite state set and one read-only head on a finite word
over an alphabet $\Sigma$, bracketed by two distinct endmarkers outside
$\Sigma$. Its head starts at the left endmarker in a designated initial
state. A transition changes the state and moves the head left, right, or
not at all; the head cannot cross an endmarker. A transition depends only
on the current state and scanned symbol. Deterministic transition rules
may be partial. A nondeterministic rule permits several successors.
There are no additional heads, writable cells, work tapes, advice, or
randomness. All states, including initial and accepting states, are counted.

Acceptance means that a finite computation reaches an accepting state;
a nonaccepting infinite computation rejects. We initially count an
accepting initial configuration as success. If success instead requires
a positive transition into an accepting state, the normalization in
Section~\ref{sec:machines} uses one additional initial copy. The main
bound below includes this extra state and holds with either convention.

For a finite set $H$, let $\R{H}$ be the monoid of all binary relations
on $H$. Products are in path order from left to right:
\[
 (x,z)\in AB
 \quad\Longleftrightarrow\quad
 \text{there exists }y\in H\text{ with }(x,y)\in A,\ (y,z)\in B.
\]
Write $\Id{F}=\{(x,x):x\in F\}$ for $F\subseteq H$; the identity is
$\Id{H}$ and the zero is the empty relation $0$. For $H=[h]$ define
\begin{equation}
\label{eq:liveness-language}
 \OWL_h=\{R_1\cdots R_\ell\in(\R{[h]})^*:
                  R_1\cdots R_\ell\ne0\}.
\end{equation}
The product of the empty word is $\Id{[h]}$, so $\eps\in\OWL_h$.

\begin{theorem}
\label{thm:main}
For every integer $h\ge2$, the language $\OWL_h$ is recognized by a
nondeterministic automaton with $h+3$ states using no left moves. Every
equivalent deterministic two-way automaton with $s$ states satisfies
\begin{equation}
\label{eq:main-bound}
 4(s+2)^2\ge 2^{\lfloor(h-2)/31\rfloor}.
\end{equation}
If an initially accepting configuration counts as acceptance, $s+2$ in
\eqref{eq:main-bound} may be replaced by $s+1$.
\end{theorem}

\begin{corollary}
\label{cor:no-polynomial}
There are no absolute constants $C,c>0$ such that every $n$-state two-way
nondeterministic automaton over every finite alphabet has an equivalent
two-way deterministic automaton with at most $Cn^c$ states.
\end{corollary}

\begin{proof}
Put $n=h+3$ in Theorem~\ref{thm:main}. Every deterministic equivalent
for this family has at least
\[
 2^{\frac12\lfloor(n-5)/31\rfloor-1}-2
\]
states. This bound exceeds $Cn^c$ for every fixed $C,c$ when $n$ is large
enough. For each $h$ the alphabet has $2^{h^2}$ symbols and is finite.
Since the constants in the proposed simulation must be independent of
the alphabet, the family is within its scope.
\end{proof}

